\documentclass[11pt]{article}
\usepackage{setspace}
\usepackage{geometry}
\geometry{margin=1in}
\setstretch{1.15}

\begin{document}

% --- Applicant Information ---
\noindent
\textbf{Zhongli Wang} \\
Phone: +49 160 9541 6015 \\
Email: nkwangzhongli@gmail.com \\
Website: https://sites.google.com/view/zhongli-wang \\
Bielefeld, Germany \\

\vspace{0.45cm}

% --- Recipient Information ---
\noindent
\textbf{Associate Professor Meike Janssen} \\
Consumer \& Behavioural Insights Group (CBIG) \\
Department of Management, Society and Communication \\
Copenhagen Business School \\
Copenhagen, Denmark \\

\vspace{0.35cm}

% --- Subject Line ---
\noindent
\textbf{Application for Postdoc (TIPRO) --- Experimental Consumer Research on Sustainable Food Consumption} \\

\vspace{0.25cm}

\noindent Dear Professor Janssen and Members of the Recruitment Committee,\\[0.2cm]

I am writing to apply for the Postdoc position in experimental consumer research within the TIPRO project at Copenhagen Business School. I am currently completing my Ph.D.\ in Economics at Bielefeld University and Universit\`a Cattolica del Sacro Cuore (Marie Sk{\l}odowska--Curie EPOC programme). My work combines behavioural economics, applied microeconometrics, and computational modelling to study how beliefs, incentives, and choice environments shape individual decision-making. I am excited about TIPRO's aim to design and test interventions for accelerating sustainable food consumption in collaboration with academic and food-sector partners across Denmark, Belgium, and the Netherlands.

\vspace{0.15cm}
My main strengths and abilities are: (i) strong quantitative training in applied microeconometrics and behavioural methods; (ii) solid experimental design literacy (randomization, treatment design, measurement of beliefs and preferences, power and robustness thinking) and hands-on experience with behavioural data analysis; (iii) experience building reproducible data pipelines and managing complex projects in collaborative settings; and (iv) a demonstrated ability to connect behavioural mechanisms to policy-relevant interventions and publishable research questions.

\bigskip

My formal training includes Experimental Economics and quantitative behavioural research. During my Master’s studies, I wrote a semester paper on the relationship between overconfidence/underconfidence and procrastination, documenting a strong link between biased self-beliefs and time-inconsistent behaviour. This early work shaped my interest in heterogeneity in psychological mechanisms---why interventions work for some groups but fail for others---which directly aligns with TIPRO’s focus on changing consumer behaviour within and across food environments. In my Ph.D., while my primary empirical setting is technology-driven markets, my research centres on behavioural mechanisms that translate naturally to sustainable consumption: belief formation, bounded rationality, social influence, and adaptive decision rules under complex choice environments.

\bigskip

My dissertation, \textit{Strategic Behavior and Economic Dynamics in ICT Markets: An Agent-Based Framework}, develops models in which heterogeneous agents update beliefs and make choices under information frictions and constraints. Methodologically, I combine simulation-based approaches with empirical validation and modern statistical learning tools to identify behavioural parameters and to evaluate counterfactual policy scenarios. This background is useful for TIPRO in two ways: first, it provides a disciplined framework for formulating behavioural hypotheses and translating them into experimentally testable interventions; second, it equips me to analyse heterogeneous treatment effects and mechanisms across consumer segments and contexts.

\bigskip

In addition to my dissertation, I have coauthored two projects that demonstrate my ability to work with large datasets, modern quantitative tools, and reproducible workflows. \textit{Polytope Fraud Theory} applies supervised machine learning to detect structured anomalies in firm financial reports, reflecting strong skills in data processing, feature engineering, and interpretable modelling. A second project develops \textit{agentic workflows} that integrate large-language-model agents into empirical research pipelines, including automated coding, text classification, and documentation. These experiences translate well to TIPRO’s collaborative setting: managing multi-source data, documenting analysis decisions transparently, and delivering policy-relevant outputs that can be communicated to academic and non-academic partners.

\bigskip

If selected, I would be eager to contribute to TIPRO by (i) co-developing intervention designs grounded in behavioural theory (e.g., choice architecture, information provision, social norms, defaults, friction/reducing barriers), (ii) implementing and analysing consumer experiments across food environments, (iii) quantifying heterogeneous effects across socio-demographic groups and contexts with credible inference strategies, and (iv) translating results into publishable manuscripts and actionable insights for partners. I am comfortable working in interdisciplinary teams and enjoy the combination of rigorous research and real-world relevance that characterizes the Consumer \& Behavioural Insights Group and the broader TIPRO consortium.

\bigskip

Thank you so much for considering my application! I would welcome the opportunity to discuss how my quantitative and behavioural research background can support TIPRO’s experimental programme and contribute to high-quality publications at the intersection of consumer behaviour, food policy, and sustainability.

\vspace{6mm}

\noindent Sincerely,\\[12mm]
\noindent Zhongli Wang

\end{document}